\section{Ablations and causal mechanism claims}\label{sec:C26:ablations}
An ablation removes or changes a mechanism, but it can also change resources, information, and optimization opportunities. Removing a retrieval module may shorten the context and reduce cost; removing a critic may reduce model calls; replacing a hypergraph with flat retrieval may alter both representation and access policy. The observed difference identifies the combined treatment unless additional controls isolate the components.

HyperSkill provides a useful example. Its no-hypergraph condition replaces the structured retrieval pipeline with flat embedding retrieval, so the result does not isolate hyperedge representation alone~\cite{xu2026hyperskill}. A stronger comparison keeps decomposition, candidate pools, ranking budget, and maintenance cost comparable while varying the storage relation. This does not invalidate the system result; it narrows the causal explanation supported by that ablation.

Oracle controls estimate headroom. Oracle retrieval tests whether downstream reasoning can use sufficient evidence; oracle state tests whether partial observation is the bottleneck; an oracle selector tests candidate-generation potential. These controls are deliberately privileged. Their performance cannot be described as deployable capability, but the gaps help prioritize engineering work. A large oracle-retrieval gap suggests retrieval deserves attention before increasing reasoning depth.

Negative controls test whether a supposed mechanism is doing meaningful work. Shuffle memory, provide irrelevant skills, randomize delegation messages, or replace a feature description with a plausible wrong label. Some controls can damage the system in several ways, so interpret them carefully. The aim is not to produce a dramatic drop but to distinguish a specific explanation from alternatives such as extra context, extra compute, or evaluator bias.

A mechanism claim requires explicit operating conditions and failure regimes. It identifies where the mechanism helps, where a simpler baseline matches it, and where it causes regressions. This is more useful than a universal superiority claim because deployments occupy particular workloads. Method comparisons should therefore preserve heterogeneity across horizon, task structure, evidence quality, and tool reliability rather than compress everything into one average.

